\begin{lemma}[Lemma 3.1]
Let $n\ge2$, $p>1$, $\varepsilon\ge0$, $0<\gamma_0\le\gamma_1$ and
$0\le\beta<1$, and suppose the Sobolev--Poincar\'e inequality on balls holds in
dimension $n$ (\noderef{CitedSobolevPoincareBall}).  Let $\Omega\subseteq\R^n$
be open, let
$A\colon\R^n\to\R^n$ satisfy (Q1) of the thesis--(Q2) of the thesis with these parameters
(\noderef{EllipticStructureQ12}) and $A(0)=0$, and assume the cited existence
and regularity result (C3) (\noderef{CitedDirichletHolder}).  Let
$\overline B_R(x_0)\subseteq\Omega$ with $R>0$ and let
$f\in C^\beta(\partial B_R(x_0))$ (\noderef{MemHolderOn}).  Then the problem
\[
\begin{cases}
\operatorname{div}A(\nabla u)=0 & \text{in }B_R,\\
u=f & \text{on }\partial B_R
\end{cases}
\]
has a solution $u$, unique in $B_R$, with
$u\in W^{1,p}_{\mathrm{loc}}(B_R)\cap C^1(B_R)$
(\noderef{IsWeakSolutionOn}, \noderef{MemW1pLocOn}) whose continuous extension
$U$ to $\overline B_R$ lies in $C^\beta(\overline B_R)$ and satisfies $U=f$ on
$\partial B_R$.  Moreover
\[
\operatorname*{osc}(u,B_R)\;\le\;\operatorname*{osc}(f,\partial B_R).
\]
Uniqueness means: any $u'$ satisfying the same list of conditions with the same
boundary datum $f$ agrees with $u$ on $B_R$.  The boundary condition is encoded,
as everywhere in this development, by the pointwise agreement of the continuous
extension with $f$ on the sphere ($\partial B_R$ is Lebesgue-null, so an
almost-everywhere condition there would be vacuous).  The oscillation bound is
written out as
$U(x)-U(y)\le\sup_{\partial B_R}f-\inf_{\partial B_R}f$ for all
$x,y\in B_R$, which is $\operatorname*{osc}(u,B_R)\le\operatorname*{osc}(f,\partial B_R)$;
$f$ is continuous on the compact sphere, so its supremum and infimum there are
finite.

The hypothesis \noderef{CitedSobolevPoincareBall} is the cited result (C6) of the
paper's Note for the formalization, carried here as an explicit argument because
the proof uses it: it enters through \noderef{BallComparisonPrinciple}, whose
statement requires it and which this proof invokes for the oscillation bound and
for uniqueness.  It is a statement about $n$ alone, and $n$ is bound first, so it
introduces no dependence on any quantity outside the list the paper names and
carries no numerical content of its own.
\end{lemma}
\begin{proof}
Write $B_r=B_r(x_0)$, $S_r=\partial B_r(x_0)$ and $\overline B_r$ for the closed
ball, and set
\[
M_0=\sup_{S_R}f,\qquad m_0=\inf_{S_R}f .
\]
By \noderef{MemHolderOn} the datum $f$ is continuous on $S_R$, and $S_R$ is
compact and nonempty (as $R>0$ and $n\ge2$), so $M_0$ and $m_0$ are finite and
attained.

\emph{Step 1: existence and regularity, from (C3).}  Apply
\noderef{CitedDirichletHolder} with the parameters $n,p,\varepsilon,\gamma_0,
\gamma_1$ and the field $A$.  Its hypotheses are exactly the ones we have:
$A$ satisfies (Q1) of the thesis--(Q2) of the thesis on $\R^n$
(\noderef{EllipticStructureQ12}), $\varepsilon\ge0$, $A(0)=0$, $R>0$,
$0\le\beta<1$ and $f\in C^\beta(S_R)$ (\noderef{MemHolderOn}).  It produces
$u$, $U$ and $g_u$ with
\begin{enumerate}
\item $u\in W^{1,p}_{\mathrm{loc}}(B_R)$ with weak gradient $g_u$
  (\noderef{MemW1pLocOn});
\item $u$ a weak solution of $\operatorname{div}A(\nabla u)=0$ in $B_R$
  (\noderef{IsWeakSolutionOn});
\item $u\in C^1(B_R)$;
\item $U\in C^\beta(\overline B_R)$ (\noderef{MemHolderOn}), $U=u$ on $B_R$, and
  $U=f$ on $S_R$.
\end{enumerate}
These are the first six assertions of the lemma verbatim.  It remains to prove
the oscillation bound and the uniqueness clause, which (C3) does not supply and
which the paper's Note explicitly requires to be proved.  Both come from
\noderef{BallComparisonPrinciple}, and both need the following device, because
that node requires its two compared functions to lie in $W^{1,p}$ of the ball it
is applied to, whereas (1) gives $u$ only in $W^{1,p}_{\mathrm{loc}}(B_R)$.

\emph{Step 2: the exhaustion device.}  Fix $0<r<R$.  We record four facts; each
is an instance of one of the ambient closure properties of the Sobolev
predicates, applied to the present data.  Note throughout that $p\ge1$, that
$B_r$ and $B_R$ are open, that $B_r\subseteq B_R$, and that
$\overline B_r$ is compact with $\overline B_r\subseteq B_R$ because $r<R$.

\emph{(2a) $u\in W^{1,p}(B_r)$ with weak gradient $g_u$, and
$u\in W^{1,p}_{\mathrm{loc}}(B_r)$ with weak gradient $g_u$.}  Apply
\noderef{WeakGradientRestrict} with $n$, $q=p$, $\Omega=B_R$, $\Omega'=B_r$ and
the pair $(u,g_u)$.  Its hypotheses are $1\le n$, $1\le p$, openness of $B_R$
and of $B_r$, and $B_r\subseteq B_R$, all of which hold.  Its clause (4) has the
further hypotheses that $\overline{B_r}$ be compact and contained in $B_R$,
which hold as just noted, and the hypothesis
$u\in W^{1,p}_{\mathrm{loc}}(B_R)$ with weak gradient $g_u$, which is (1) of
Step 1; it concludes $u\in W^{1,p}(B_r)$ with weak gradient $g_u$
(\noderef{MemW1pOn}).  Its clause (2), from the same hypothesis, concludes
$u\in W^{1,p}_{\mathrm{loc}}(B_r)$ with weak gradient $g_u$
(\noderef{MemW1pLocOn}); this is the first clause of
\noderef{IsWeakSolutionOn} on $B_r$, needed in (2b).

\emph{(2b) $u$ is a weak solution in $B_r$.}  The first clause of
\noderef{IsWeakSolutionOn} on $B_r$ is the second conclusion of (2a).  For the
integral identity, let $\varphi\in W^{1,p}(B_r)$ with weak gradient $g_\varphi$
(\noderef{MemW1pOn}) have compact support with
$\operatorname{supp}\varphi\subseteq B_r$; these are exactly the conditions the
test class of \noderef{IsWeakSolutionOn} on $B_r$ imposes.  Apply
\noderef{WeakGradientZeroExtend} with $n$, $q=p$, $\Omega=B_R$, the ball $B_r$
and the pair $(\varphi,g_\varphi)$: its hypotheses are $1\le n$, $1\le p$,
openness of $B_R$, $B_r\subseteq B_R$, $\varphi\in W^{1,p}(B_r)$ with weak
gradient $g_\varphi$, and the compact support condition, all of which hold.  It
gives
\[
\varphi\in W^{1,p}(B_R)\quad\text{with weak gradient}\quad
G=\mathbf1_{B_r}g_\varphi ,
\]
and $\varphi$ has compact support with
$\operatorname{supp}\varphi\subseteq B_r\subseteq B_R$.  So the pair
$(\varphi,G)$ lies in the test class of \noderef{IsWeakSolutionOn} on $B_R$,
and the identity for $u$ on $B_R$ --- which is the hypothesis (2) of Step 1 ---
gives
\[
\int_{B_R}\Bigl[\sum_{i=1}^n A^i\bigl(g_u(x)\bigr)\,G^i(x)+0\cdot\varphi(x)\Bigr]dx=0 .
\]
The integrand vanishes off $B_r$: there $G=0$ by construction and $\varphi=0$
because $\operatorname{supp}\varphi\subseteq B_r$.  An integral over a set
depends on the integrand only through its product with that set's indicator, and
those products for $B_R$ and for $B_r$ coincide here, so the integral over $B_R$
equals the integral over $B_r$; and on $B_r$ we have $G=g_\varphi$.  Hence
\[
\int_{B_r}\Bigl[\sum_{i=1}^n A^i\bigl(g_u(x)\bigr)\,g_\varphi^i(x)+0\cdot\varphi(x)\Bigr]dx=0 ,
\]
which is the integral identity of \noderef{IsWeakSolutionOn} on $B_r$ for the
test pair $(\varphi,g_\varphi)$.

\emph{The integrability clause that \noderef{IsWeakSolutionOn} also carries.}
That node asks, for each admissible test pair, that the integrand be Lebesgue
integrable on the domain before asserting the identity, so that the identity
cannot be satisfied by a junk value.  Here it holds, and the $C^1$ regularity
already in hand is what gives it.  Since $B\equiv0$ the integrand is
$x\mapsto\sum_iA^i(g_u(x))g_\varphi^i(x)$.  It vanishes almost everywhere off
$\operatorname{supp}\varphi$: the set $B_r\setminus\operatorname{supp}\varphi$ is
open and $\varphi$ vanishes identically on it, so its weak gradient
$g_\varphi$ vanishes almost everywhere there (Preamble: du
Bois--Reymond, second recorded consequence).  For the bound on $\operatorname{supp}\varphi$ the route has to be
spelled out, because the $C^1$ datum of Step 1 is an assertion about $u$ and not
about $g_u$.  Write $\nabla u$ for the \emph{classical} gradient of $u$ on
$B_R$, so that $(\nabla u)^i(x)=D_iu(x)$ is the derivative of $u$ at $x$ in the
direction of the $i$-th standard basis vector.

\emph{First, $\nabla u$ is a weak gradient of $u$ on $B_R$.}  Apply
\noderef{WeakGradientOfC1} with the present $n$, with $q=p$, with the open set
$B_R$, with the function $u$ and with the field $\nabla u$.  Its hypotheses
hold: $1\le n$ and $1\le p$, from $2\le n$ and $1<p$; $B_R$ is open; $u$ is
differentiable at every point of $B_R$ with $D_iu(x)=(\nabla u)^i(x)$ for every
such $x$ and every $i$, which is assertion (3) of Step 1, $u\in C^1(B_R)$, read
on the open set $B_R$; and $\nabla u$ is continuous on $B_R$, which is the other
half of that same assertion.  Its conclusion is
$u\in W^{1,p}_{\mathrm{loc}}(B_R)$ with weak gradient $\nabla u$
(\noderef{MemW1pLocOn}).

\emph{Second, $g_u$ and $\nabla u$ agree almost everywhere.}  Assertion (1) of
Step 1 gives $u\in W^{1,p}_{\mathrm{loc}}(B_R)$ with weak gradient $g_u$, and
the previous paragraph gives the same with $\nabla u$ in its place; both are
therefore weak gradients of the one function $u$ on the one open set $B_R$.  A
weak gradient is determined almost everywhere on an open set (Preamble: du
Bois--Reymond, second recorded consequence, whose first sentence is exactly
this), so
\[
g_u=\nabla u\qquad\text{almost everywhere on }B_R .
\tag{2.1}
\]
\emph{This is as much as can be had, and it is all that is needed.}  The field
$g_u$ is a weak gradient, alterable on a null set without disturbing any
hypothesis of this lemma, so $g_u$ is \emph{not} continuous and assertion (3)
does not make it so; what (2.1) supplies is that $g_u$ agrees almost everywhere
with a field continuous on $B_R$, hence that $g_u$ is \emph{essentially}
bounded on compact subsets of $B_R$.  The same argument applies verbatim to any
pair $(u^\flat,g_{u^\flat})$ satisfying assertions (1) and (3), which is how
Step 6 uses it for $u'$.

Now $\nabla u$ is continuous on $B_R$ and
$\operatorname{supp}\varphi\subseteq B_r\subseteq B_R$ is compact, so $\nabla u$
is bounded on $\operatorname{supp}\varphi$ and its image there lies in a compact
subset $Q\subseteq\R^n$; $A$ is continuous on $\R^n$ by clause \texttt{hcont} of
\noderef{EllipticStructureQ12}, so $|A|$ is bounded on $Q$, say by $M$.  With
$(2.1)$ this gives $|A(g_u(x))|\le M$ for almost every
$x\in\operatorname{supp}\varphi$.  And $x\mapsto A(g_u(x))$ is measurable, by
that node's Carath\'eodory clauses together with the measurability of $g_u$;
that half of the argument never needed the $C^1$ datum.

Finally the domination.  The integrand
$x\mapsto\sum_iA^i(g_u(x))g_\varphi^i(x)$ is measurable on $B_r$, being a finite
sum of products of measurable functions, and almost everywhere on $B_r$ it
satisfies
\[
\Bigl|\sum_iA^i(g_u(x))g_\varphi^i(x)\Bigr|\ \le\ M\,|g_\varphi(x)| :
\]
almost everywhere on $\operatorname{supp}\varphi$ by the bound just obtained
together with the Cauchy--Schwarz inequality (Preamble), and almost everywhere
on $B_r\setminus\operatorname{supp}\varphi$ because $g_\varphi$ vanishes almost
everywhere there, as recorded above, so that both sides vanish.  The dominating
function is integrable: $g_\varphi\in L^p(B_r)$ and $|B_r|<\infty$ give
$g_\varphi\in L^1(B_r)$ by H\"older's inequality on a set of finite measure
(Preamble).  So the integrand is integrable on $B_r$ by the Preamble's
domination clause~(a), which asks for precisely an almost-everywhere bound by an
integrable function.  As
the test pair was arbitrary, $u$ is a weak
solution of $\operatorname{div}A(\nabla u)=0$ in $B_r$
(\noderef{IsWeakSolutionOn}).

\emph{(2c) Every constant is a weak solution in $B_r$, and this is where
$A(0)=0$ is used.}  Let $c\in\R$ and let $\kappa$ denote the constant function
with value $c$.  The ball $B_r$ is bounded, so $|B_r|<\infty$.  Apply
\noderef{WeakGradientLinear} with $n$, $q=p$, $\Omega=B_r$ and this $c$: its
clause (4), whose extra hypothesis $|B_r|<\infty$ holds, gives
$\kappa\in W^{1,p}(B_r)$ with weak gradient $0$ (\noderef{MemW1pOn}), and its
clause (3) gives $\kappa\in W^{1,p}_{\mathrm{loc}}(B_r)$ with weak gradient $0$
(\noderef{MemW1pLocOn}), which is the first clause of
\noderef{IsWeakSolutionOn}.  For the equation, the integral identity of
\noderef{IsWeakSolutionOn} for $\kappa$ with weak gradient $0$ reads
\[
\int_{B_r}\Bigl[\sum_{i=1}^nA^i(0)\,g_\varphi^i+0\cdot\varphi\Bigr]dx=0 ,
\]
which holds because $A(0)=0$ makes the integrand identically zero; and the
integrability clause that \noderef{IsWeakSolutionOn} also carries holds for the
same reason, the zero function being integrable on $B_r$.  This is the
only place in the \emph{comparison} argument where $A(0)=0$ is consumed --- it is
also a hypothesis of \noderef{CitedDirichletHolder}, used in Step 1 --- and it is
\S3's own normalisation, available because $\operatorname{div}A(0)=0$ for any
constant field.

\emph{(2d) Adding a constant to a solution gives a solution.}  Let $\tilde u$ be
a weak solution in $B_r$ with weak gradient $g$ (\noderef{IsWeakSolutionOn}) and
let $c\in\R$, with $\kappa$ as in (2c).  Apply \noderef{WeakGradientLinear} with
$n$, $q=p$, $\Omega=B_r$, $a=b=1$, the pair $(f_1,g_1)=(\tilde u,g)$ and the
pair $(f_2,g_2)=(\kappa,0)$.  Its clause (1) has as hypotheses
$\tilde u\in W^{1,p}_{\mathrm{loc}}(B_r)$ with weak gradient $g$ --- the first
clause of \noderef{IsWeakSolutionOn} for $\tilde u$ --- and
$\kappa\in W^{1,p}_{\mathrm{loc}}(B_r)$ with weak gradient $0$, which is clause
(3) of the same node as recorded in (2c); it concludes that
$\tilde u+c\in W^{1,p}_{\mathrm{loc}}(B_r)$ with weak gradient
$1\cdot g+1\cdot0=g$.  If moreover $\tilde u\in W^{1,p}(B_r)$ with weak gradient
$g$, clause (2) of the same node, with the $W^{1,p}$ membership of $\kappa$ from
(2c), gives $\tilde u+c\in W^{1,p}(B_r)$ with weak gradient $g$.  For the
equation, the integral identity of \noderef{IsWeakSolutionOn} involves the
function only through $A(x,\tilde u(x),g(x))=A(g(x))$ and through $B\equiv0$;
since the operator of \S3 does not depend on the function value, and since
$\tilde u+c$ has the same weak gradient $g$, the integrand for $\tilde u+c$ is
literally the same \emph{function} as the integrand for $\tilde u$.  Both
clauses that \noderef{IsWeakSolutionOn} binds at a test pair therefore transfer
at once: the integrability of that integrand on $B_r$, and the vanishing of its
integral.  So $\tilde u+c$ is a weak solution in
$B_r$ with weak gradient $g$.  (Note that this step does \emph{not} need
$A(0)=0$.)

\emph{Step 3: the upper bound $u\le M_0$ a.e.\ in $B_R$.}  Let $\eta>0$.  Since
$U$ is continuous on the compact set $\overline B_R$ it is uniformly continuous
there, so there is $\delta>0$ such that
\[
x,y\in\overline B_R,\ |x-y|<\delta\ \Longrightarrow\ |U(x)-U(y)|<\eta .
\tag{3.1}
\]
Fix any $r$ with $\max\{0,R-\delta\}<r<R$ and put
\[
M_r=\max_{S_r}U ,
\]
which exists because $U$ is continuous on $\overline B_R\supseteq S_r$ and $S_r$
is compact and nonempty.  We claim $M_r\le M_0+\eta$.  Let $z\in S_r$ and let
$z'=x_0+\frac Rr(z-x_0)$ be its radial projection onto $S_R$; then
$|z'-x_0|=R$, so $z'\in S_R$, and
\[
|z-z'|=\Bigl|\frac Rr-1\Bigr|\,|z-x_0|=\frac{R-r}{r}\cdot r=R-r<\delta .
\]
Hence by (3.1) and $U=f$ on $S_R$,
\[
U(z)\le U(z')+\eta=f(z')+\eta\le M_0+\eta ,
\]
and taking the maximum over $z\in S_r$ gives $M_r\le M_0+\eta$ as claimed.

Now apply \noderef{BallComparisonPrinciple} on the ball $B_r$ to the pair
$(u,\kappa)$, where $\kappa$ is the constant function with value $M_r$.  Its
hypotheses hold: $n\ge2$, $p>1$, $\varepsilon\ge0$, $0<\gamma_0\le\gamma_1$, the
Sobolev--Poincar\'e hypothesis \noderef{CitedSobolevPoincareBall} in dimension
$n$ --- which that node's statement requires and which is a hypothesis of the
present lemma, for the same dimension $n$ --- and
\noderef{EllipticStructureQ12} for $A$ on $\R^n$ are hypotheses of the present
lemma; $r>0$; $u$ and $\kappa$ are weak solutions in $B_r$ by (2b) and (2c) and
lie in $W^{1,p}(B_r)$ by (2a) and (2c); the continuous extensions to
$\overline B_r$ may be taken to be $U$ --- continuous on
$\overline B_r\subseteq\overline B_R$ and equal to $u$ on $B_r$ by (4) of Step 1
--- and the constant $M_r$ itself, which is continuous and extends $\kappa$; and
on $S_r$ we have $U\le M_r$ by the definition of $M_r$.  The conclusion is
\[
u\le M_r\le M_0+\eta\qquad\text{a.e.\ in }B_r .
\]
This holds for every $r\in(\max\{0,R-\delta\},R)$.  Choose a sequence
$r_j\uparrow R$ inside that interval; then $B_R=\bigcup_jB_{r_j}$, and for each
$j$ the inequality $u\le M_0+\eta$ holds off a null subset $N_j$ of $B_{r_j}$.
Since $\bigcup_jN_j$ is null, $u\le M_0+\eta$ a.e.\ in $B_R$.

Finally let $\eta=1/i$ for $i=1,2,\dots$: for each $i$ we get a null set outside
which $u\le M_0+1/i$, and the union of these null sets is null, so
\[
u\le M_0\qquad\text{a.e.\ in }B_R .
\tag{3.2}
\]

\emph{Step 4: the lower bound $u\ge m_0$ a.e.\ in $B_R$.}  Symmetrically, let
$\eta>0$, take $\delta$ from (3.1) and $r\in(\max\{0,R-\delta\},R)$, and put
$m_r=\min_{S_r}U$.  For $z\in S_r$ with radial projection $z'\in S_R$ we get
$U(z)\ge U(z')-\eta=f(z')-\eta\ge m_0-\eta$, hence $m_r\ge m_0-\eta$.  Apply
\noderef{BallComparisonPrinciple} on $B_r$ to the pair $(\kappa',u)$, where
$\kappa'$ is the constant $m_r$: the hypotheses are verified exactly as in Step 3
with the two roles exchanged --- in particular its Sobolev--Poincar\'e hypothesis
\noderef{CitedSobolevPoincareBall} is again supplied by the present lemma's
hypothesis of the same name --- the boundary comparison being $m_r\le U$ on $S_r$
by the definition of $m_r$.  The conclusion is $m_r\le u$ a.e.\ in $B_r$, so
$u\ge m_0-\eta$ a.e.\ in $B_r$, and exhausting in $r$ and then letting
$\eta\downarrow0$ exactly as in Step 3 gives
\[
u\ge m_0\qquad\text{a.e.\ in }B_R .
\tag{3.3}
\]

\emph{Step 5: the oscillation bound.}  We must show
$U(x)-U(y)\le M_0-m_0$ for all $x,y\in B_R$.  By (3.2) and $U=u$ on $B_R$ we
have $U\le M_0$ almost everywhere on $B_R$.  The function $U$ is continuous on
$B_R$ and $B_R$ is open, so the inequality upgrades to every point: if
$U(x_1)>M_0$ for some $x_1\in B_R$, then by continuity $U>M_0$ on a ball around
$x_1$ contained in $B_R$, a set of positive measure, contradicting
$U\le M_0$ a.e.  Hence $U\le M_0$ everywhere on $B_R$, and in the same way (3.3)
gives $U\ge m_0$ everywhere on $B_R$.  Therefore for all $x,y\in B_R$
\[
U(x)-U(y)\le M_0-m_0
=\sup_{\partial B_R}f-\inf_{\partial B_R}f ,
\]
which is the seventh assertion, i.e.\
$\operatorname*{osc}(u,B_R)\le\operatorname*{osc}(f,\partial B_R)$.

\emph{Step 6: uniqueness.}  Let $u'$, $U'$, $g_{u'}$ satisfy the same six
conditions with the same datum $f$: $u'$ a weak solution in $B_R$
(\noderef{IsWeakSolutionOn}), $u'\in W^{1,p}_{\mathrm{loc}}(B_R)$ with weak
gradient $g_{u'}$ (\noderef{MemW1pLocOn}), $u'\in C^1(B_R)$,
$U'\in C^\beta(\overline B_R)$ (\noderef{MemHolderOn}), $U'=u'$ on $B_R$ and
$U'=f$ on $S_R$.  We show $u'=u$ pointwise on $B_R$.

Let $\eta>0$.  Both $U$ and $U'$ are continuous on the compact $\overline B_R$,
hence uniformly continuous, so there is $\delta>0$ such that (3.1) holds for $U$
and the same statement holds for $U'$.  Fix
$r\in(\max\{0,R-\delta\},R)$.  For $z\in S_r$ with radial projection
$z'\in S_R$ we have $|z-z'|=R-r<\delta$, and $U(z')=f(z')=U'(z')$, so
\[
U(z)-U'(z)=\bigl(U(z)-U(z')\bigr)+\bigl(U'(z')-U'(z)\bigr)<\eta+\eta=2\eta ,
\]
that is
\[
U\le U'+2\eta\qquad\text{on }S_r .
\tag{3.4}
\]
Now apply \noderef{BallComparisonPrinciple} on $B_r$ to the pair
$(u,\;u'+2\eta)$.  Its hypotheses hold: $u$ is a weak solution in $B_r$ lying in
$W^{1,p}(B_r)$ by (2a) and (2b); $u'$ likewise, by the same two arguments
applied to $u'$, which satisfies the same hypotheses --- including the $C^1$
regularity that (2b)'s integrability step consumes, so that $(2.1)$ holds for
the pair $(u',g_{u'})$ by the verbatim repetition noted there; hence $u'+2\eta$ is a
weak solution in $B_r$ with weak gradient $g_{u'}$ by (2d), and it lies in
$W^{1,p}(B_r)$ because $u'$ does and constants do by (2c); the continuous
extensions to $\overline B_r$ are $U$ and $U'+2\eta$, which extend $u$ and
$u'+2\eta$ from $B_r$; its Sobolev--Poincar\'e hypothesis
\noderef{CitedSobolevPoincareBall} in dimension $n$ is the present lemma's
hypothesis of the same name; and the boundary comparison on $S_r$ is (3.4).  The
conclusion is
\[
u\le u'+2\eta\qquad\text{a.e.\ in }B_r .
\]
Exhausting in $r$ along a sequence $r_j\uparrow R$ and then letting
$\eta=1/i\downarrow0$, exactly as at the end of Step 3, gives $u\le u'$ a.e.\ in
$B_R$.  Exchanging the roles of $u$ and $u'$ --- the hypotheses are symmetric in
them, and (3.4) holds with $U$ and $U'$ exchanged by the same computation ---
gives $u'\le u$ a.e.\ in $B_R$.  Hence $u=u'$ almost everywhere in $B_R$.

Finally, $u$ and $u'$ are both continuous on $B_R$, being $C^1$ there, and two
continuous functions that agree almost everywhere on an open set agree
everywhere on it: if $u(x_1)\ne u'(x_1)$ for some $x_1\in B_R$ then
$u\ne u'$ on a ball around $x_1$ inside $B_R$, a set of positive measure.
Therefore $u'=u$ on $B_R$, which is the uniqueness assertion.
\end{proof}
